\documentclass[11pt]{ctexart}
\usepackage{amsmath,amssymb,longtable}
\begin{document}
\section*{P24: 当前完整数学模型（2026-10-20）}
\section*{修订依据与边界}以源文第21页流入流出文字恢复全网络守恒；补齐有限采购/运输能力及预测器训练域。本地营养配送实例给出有限业务边界。




\section*{保留的核心模型}
第21页上的世界粮食计划署食物篮模型，将一种营养约束条件融入到混合整数优化问题中。 \emph{在}之后可以直接求解。预测函数$\hat h(\boldsymbol{x})$可以通过任何经过训练的、能够表示MIO的学习者来实现，这些学习者可以由本文所支持的模型来提供支持（例如，线性模型、树状模型、森林模型或ReLU网络）。

\paragraph{决策变量。}
$x_k$表示食物物品$k$的数量；$F_{ijk}$是供应/运输网络上的商品流动变量；而$y$则是由嵌入式模型$\hat h(\boldsymbol{x})$生成的预测适口性评分（或预测的约束输出）。

\paragraph{集合与参数。}
$\mathcal N_{\mathcal S}$、$\mathcal N_{\mathcal T}$和$\mathcal N_{\mathcal D}$分别是来源节点、转运节点和需求节点集合。$\mathcal K$是商品集合，而$\mathcal L$是营养元素集合。成本$p_{ik}^P$和$p_{ijk}^T$分别代表了采购成本和运输成本，$\gamma$是一个转换因子，$d_i$代表每日需求量，$days$表示规划的天数，$Nutval_{kl}$表示食物$k$中营养元素$l$的含量，$Nutreq_l$表示所需的营养元素$l$的数量，而$t$则是最低可接受的学习分数阈值。为了使保留模型完全可解，我们采用了一个经过训练的ReLU网络来实现学习映射$\hat h(\boldsymbol{x})$，该网络是本文明确支持的MIO可表示学习器之一。让$\mathcal R_\ell$代表隐藏层$\ell=1,\dots,H$的神经元集合，其权重为$W_{rq}^{(\ell)}$，偏置为$b_r^{(\ell)}$，并且具有有效的预激活界限$L_r^{(\ell)} \le s_r^{(\ell)} \le U_r^{(\ell)}$。

\section*{应用层核心公式化}
\begin{equation}
\min_{\boldsymbol{x},\, y,\, \boldsymbol{F}}
\sum_{i \in \mathcal{N}_{\mathcal S}}
\sum_{j \in \mathcal{N}_{\mathcal T} \cup \mathcal{N}_{\mathcal D}}
\sum_{k \in \mathcal K} p_{ik}^P F_{ijk}
+
\sum_{i \in \mathcal{N}_{\mathcal S} \cup \mathcal{N}_{\mathcal T}}
\sum_{j \in \mathcal{N}_{\mathcal T} \cup \mathcal{N}_{\mathcal D}}
\sum_{k \in \mathcal K} p_{ijk}^T F_{ijk}
\tag{6a}
\end{equation}

\begin{equation}
\sum_{j \in \mathcal N_{\mathcal T}\cup\mathcal N_{\mathcal D}} F_{ijk}
=
\sum_{j \in \mathcal N_{\mathcal S}\cup\mathcal N_{\mathcal T}} F_{jik}
\quad \forall i \in \mathcal N_{\mathcal T},\ \forall k \in \mathcal K
\tag{6b}
\end{equation}

\begin{equation}
\sum_{j \in \mathcal N_{\mathcal S} \cup \mathcal N_{\mathcal T}} \gamma F_{jik}
= d_i x_k \, days
\quad \forall i \in \mathcal N_{\mathcal D},\ \forall k \in \mathcal K
\tag{6c}
\end{equation}

\begin{equation}
\sum_{k \in \mathcal K} Nutval_{kl} x_k \ge Nutreq_l
\quad \forall l \in \mathcal L
\tag{6d}
\end{equation}

\begin{equation}
x_{\mathrm{salt}} = 5
\tag{6e}
\end{equation}

\begin{equation}
x_{\mathrm{sugar}} = 20
\tag{6f}
\end{equation}

\begin{equation}
y \ge t
\tag{6g}
\end{equation}

\begin{equation}
y = \hat h(\boldsymbol{x})
\tag{6h}
\end{equation}

\begin{equation}
F_{ijk},\, x_k \ge 0
\quad \forall i,j \in \mathcal N,\ \forall k \in \mathcal K
\tag{6i}
\end{equation}

\section*{$\hat h(\boldsymbol{x})$ 的显式混合整数优化嵌入}
\begin{equation}
a_k^{(0)} = x_k
\quad \forall k \in \mathcal K
\tag{Emb-a}
\end{equation}

\begin{equation}
s_r^{(\ell)} =
\sum_{q \in \mathcal R_{\ell-1}} W_{rq}^{(\ell)} a_q^{(\ell-1)} + b_r^{(\ell)}
\quad \forall \ell=1,\dots,H,\ \forall r \in \mathcal R_\ell
\tag{Emb-b}
\end{equation}

\begin{equation}
a_r^{(\ell)} \ge s_r^{(\ell)}
\quad \forall \ell=1,\dots,H,\ \forall r \in \mathcal R_\ell
\tag{Emb-c}
\end{equation}

\begin{equation}
a_r^{(\ell)} \le s_r^{(\ell)} - L_r^{(\ell)} \left(1-\delta_r^{(\ell)}\right)
\quad \forall \ell=1,\dots,H,\ \forall r \in \mathcal R_\ell
\tag{Emb-d}
\end{equation}

\begin{equation}
a_r^{(\ell)} \le U_r^{(\ell)} \delta_r^{(\ell)}
\quad \forall \ell=1,\dots,H,\ \forall r \in \mathcal R_\ell
\tag{Emb-e}
\end{equation}

\begin{equation}
a_r^{(\ell)} \ge 0
\quad \forall \ell=1,\dots,H,\ \forall r \in \mathcal R_\ell
\tag{Emb-f}
\end{equation}

\begin{equation}
y =
\sum_{q \in \mathcal R_H} W_q^{(H+1)} a_q^{(H)} + b^{(H+1)}
\tag{Emb-g}
\end{equation}

\begin{equation}
\delta_r^{(\ell)} \in \{0,1\}
\quad \forall \ell=1,\dots,H,\ \forall r \in \mathcal R_\ell
\tag{Emb-h}
\end{equation}

\section*{解释}
在第21页上的应用级模型，一旦约束条件(6h)被上述显式预测器嵌入所替代，就会成为一个完整的独立MIO模型。该论文还支持其他可表示为MIO的学习器，但经过训练的ReLU网络是一个完全指定的保留实例，因此之前缺失的学习约束块也得到了填补。



\section*{有限业务数据域}
实例逐项给出有向运输弧、价格、食品营养、需求、采购上限$Cap^P$、弧吨位上限$Cap^T$，
食品上限$U$与类别集合$K_g$及类别上下限$L_g,U_g$。
\[
\sum_jF_{ijk}\le Cap^P_{ik}\ (i\in N_S),\quad
\sum_kF_{ijk}\le Cap^T_{ij},\quad0\le x_k\le U_k,\quad
L_g\le\sum_{k\in K_g}x_k\le U_g.
\]
这些参数给出有限采购、运输能力及学习器的有效输入域。
ReLU权重、偏置和合法预激活界由实例逐层给出；上文所有嵌入约束同时成立。

\end{document}
